\begin{abstract}
%Pre-trained language models (PTLMs) have achieved state-of-the-art performance over a range of natural language processing tasks. These models are typically trained over a static generic-domain corpus up to a certain date (\eg, a Wikipedia snapshot) and fixed for future applications. However, in practice, data from new domains or later time steps, which both feature an evolving data distribution, may constantly emerge and carry new knowledge. Modeling such emerging knowledge while retaining past knowledge is crucial to downstream task performance over all seen distributions. Towards this goal, we study the problem of lifelong learning an up-to-date language model, where the model visits a sequence of unlabeled corpora from various domains or incremental time steps. We propose two data streams, a domain-incremental \textit{academic papers} data stream and a temporally ordered \textit{social media} data stream for study. Extending from replay memory based lifelong learning approaches, we show that using \textit{dark knowledge} (\eg, attention weights, similarity between examples) of old models for future regularization can be helpful to retaining past performance. We observe positive improvements in some of our experiment setups, while we found the simple experience replay (ER) approach to be very effective as well in general despite its simplicity. We hope our work provides a solid and realistic problem setup and inspires future works to study more powerful approaches. 

Pretrained language models (PTLMs) are typically learned over a large, static corpus and further fine-tuned for various downstream tasks.
% While in an open-ended real-world setup, data from new domains or later time steps constantly emerge and carry new knowledge.
%% by Xiang
However, when deployed in the real world, a PTLM-based model must deal with data distributions that deviate from what the PTLM was initially trained on. %, or newly emerged data that contains out-of-distribution information.
In this paper, we study a \textit{lifelong language model pretraining} challenge where a PTLM is continually updated so as to adapt to emerging data. 
Over a domain-incremental research paper stream and a chronologically-ordered tweet stream, we incrementally pretrain a PTLM with different continual learning algorithms, and keep track of the downstream task performance (after fine-tuning). We evaluate PTLM's ability to adapt to new corpora while retaining learned knowledge in earlier corpora.
% The resulting model is fine-tuned over various downstream tasks to measure retention of past knowledge and acquisition of new knowledge. 
%\xiang{1 sentence to summarize the key findings and takeways.}
Our experiments show distillation-based approaches to be most effective in retaining downstream performance in earlier domains. The algorithms also improve knowledge transfer, allowing models to achieve better downstream performance over the latest data, and improve temporal generalization when distribution gaps exist between training and evaluation because of time.
% We evaluate existing continual learning algorithms, and specifically focus on distillation based approaches, which regularizes the change of ``dark knowledge'' in  models. The approaches not only achieves the best performance, but also dissect dark knowledge that are relevant to forgetting. 
We believe our problem formulation, methods, and analysis will inspire future studies towards continual pretraining of language models.
% We hope our problem formulation, methods, and analysis could inspire future study towards continual learning of language models.

%We propose two data streams, a domain-incremental \textit{academic papers} data stream and a temporally ordered \textit{social media} data stream for study. Extending from replay memory based lifelong learning approaches, we show that using \textit{dark knowledge} (\eg, attention weights, similarity between examples) of old models for future regularization can be helpful to retaining past performance. We observe positive improvements in some of our experiment setups, while we found the simple experience replay (ER) approach to be very effective as well in general despite its simplicity. We hope our work provides a solid and realistic problem setup and inspires future works to study more powerful approaches. 


\end{abstract}